% !PRINTED 137
\weberpage{137}
%<<CONT
few non-metacyclic groups contained among them, and especially the simple non-cyclic groups, therefore command all the greater interest.

It has not so far been possible to comprehend the degree numbers of the simple groups in a definite law, and one has been content to determine all simple groups up to certain limits of the degree numbers, on the one hand with the help of the theorems developed above, on the other hand by special methods. \gesperrt{Hölder} has carried out these investigations for the degree numbers up to 200, and \gesperrt{Cole} for the degree numbers up to 500\footnote{\gesperrt{Hölder}, Mathematische Annalen, Bd.\,40. \gesperrt{Cole}, American Journal, Vol.\,14.}. Of simple non-cyclic groups, only those of degrees 60, 168, 360, already known from other investigations, have emerged in this. Further, a simple group of degree 660 is also known, and recently \gesperrt{Cole} and \gesperrt{Moore} have called attention to a simple group of degree 504\footnote{Bulletin of the New York mathem. society, Oct.\,1893.}. These investigations become very laborious as the degree numbers grow. We shall confine ourselves here to the consideration of the degree numbers of the first hundred.

Here, by the three general theorems of \gesperrt{Sylow} and \gesperrt{Frobenius}, the majority of the groups prove to be metacyclic, and there remain only the degree numbers

\[ 36,\ 60,\ 72,\ 84,\ 90,\ 100 \]

But that a group of degree $36$ cannot be simple follows from \gesperrt{Sylow}'s theorem. For such a group would have to have a divisor of degree $9$, and would therefore have to be representable by the permutations of four figures. But that is impossible, because there are only 24 permutations of four figures. A group of degree $72 = 8 . 9$ must, by \S.\,30, II., since 8 is not $\equiv 1 \pmod{3}$, contain a divisor of degree 18, and would therefore, if it were simple, likewise be representable by the permutations of four figures, which is still less possible. That simple groups of degrees 84 and 100 cannot exist follows likewise at once from \gesperrt{Sylow}'s theorems, since, if we separate out the groups of degree $7$ or of degree $25$, no divisor remains that is congruent to 1 with respect to the modulus 7 or 5. That all these degree numbers can belong only to metacyclic groups then follows without further ado from \S.\,6, 2., from the fact that the groups whose degrees
%CONT>>
